\section{PBW theorem}

Let $\g$ be any Lie algebra over a field $k$.

\begin{definition}
The \emph{universal enveloping algebra} $\U\g$ of $\g$ is the associative
algebra over $k$ generated by $\g$ with relations $xy - yx = [x, y]$ for all
$x, y \in \g$.
\end{definition}

More formally, if $V$ is a vector space over $k$, then
\[
  TV = k + V + V \otimes V + V^{\otimes 3} + \dots = \bigoplus_{n \geq 0} V^{\otimes n}
\]
is the \emph{tensor algebra of $V$}, and is the free associative algebra
generated by $V$. Let $J$ be the two-sided ideal in $T\g$ generated by
$x \otimes y - y \otimes x - [x, y]$ for all $x, y \in \g$; then
$\U\g = T\g/J$.

\begin{exercise}
An \emph{enveloping algebra} for $\g$ is a linear map $i\colon \g \to A$, where
$A$ is an associative algebra, and $i(x)i(y) - i(y)i(x) = i([x, y])$. (For
example, if $V$ is a representation of $\g$, we can take $A = \End(V)$ and $i$
the action map.) Show that $\U\g$ is initial in the category of enveloping
algebras, i.e.\ any such map $i$ factors uniquely through $\U\g$.
\end{exercise}

\begin{remark}
This is old terminology: anywhere other than the above exercise and old
textbooks, ``enveloping algebra'' almost certainly means ``universal
enveloping algebra''.
\end{remark}

The Casimir operator naturally lives in $\U\g$; in fact, as we've shown, in
the centre $Z(\U\g)$.

Observe that $T\g$ is a graded algebra, but that the relations we are imposing
on it are not homogeneous ($x \otimes y$ and $y \otimes x$ have degree $2$,
while $[x, y]$ has degree $1$). Consequently, $\U\g$ is not graded. It is,
however, \emph{filtered}.

Define $(\U\g)_n$ to be the span of products of $\leq n$ elements of $\g$;
then $\U_n \subseteq \U_m$ for $n \leq m$, and
$\U_n \cdot \U_m \subseteq \U_{n+m}$. (In particular, $\U_0 = k$,
$\U_1 = k + \g$, etc.)

\begin{exercise}
Show that if $x \in \U_m$ and $y \in \U_n$, then $[x, y] \in \U_{n+m-1}$ (we
clearly have $[x, y] \in \U_{n+m}$, and are asserting that in fact the degree
is one lower). By abuse of notation, for $x, y \in \U\g$ we write $[x, y]$ to
mean $xy - yx$.
\end{exercise}

\begin{definition}
For a filtration $F_0 \subseteq F_1 \subseteq \cdots$, we write
$\gr(F) = \bigoplus F_i/F_{i-1}$. (We probably need $F_i$ to be a filtration
of an algebra for this to make sense\dots) This is the \emph{associated graded
algebra}.
\end{definition}

\begin{theorem}[Poincar\'e--Birkhoff--Witt]
\[
  \gr \U\g = \bigoplus_n (\U\g)_n/(\U\g)_{n-1} \cong S\g.
\]
Equivalently, if $\{x_1, \dots, x_n\}$ is a basis for $\g$, then
$\setc{x_1^{a_1}x_2^{a_2}\cdots x_n^{a_n}}{a_i \in \Z_{\geq 0}}$ is a basis for
$\U\g$. In particular, $\g \hookrightarrow \U\g$.
\end{theorem}

\begin{exercise}\leavevmode
\begin{enumerate}
  \item Show that the previous exercise,
    $[\U_n, \U_m] \subseteq \U_{n+m-1}$, implies that there is a well-defined
    map $S\g \to \gr \U\g$ extending the map $x \mapsto x$ on $\g$.
  \item Show that the map is surjective, or equivalently that the
    $\{\prod x_i^{a_i}\}$ span $\U\g$.
\end{enumerate}
\end{exercise}

Therefore, the hard part of the PBW theorem is injectivity. We will not prove
the theorem.

\begin{remark}
The theorem should imply that $\U\g$ is a deformation of $S\g$ (i.e.\ if we
consider $xy - yx = t[x, y]$ then for $t = 0$ we get $S\g$, whereas for
$t = 1$, and in fact for generic $t$ by renormalising, we get $\U\g$). The
shortest proof of the theorem is indeed via deformation theory.
\end{remark}

\section{Back to representations}

\begin{exercise}
If $V$ is a representation of $\g$ and $v \in V$, then the $\g$-submodule of
$V$ generated by $v$ is $\U\g \cdot v$.
\end{exercise}

\lecture{18}

\begin{definition}
$V$ is a \emph{highest weight module} for $\g$ if there exists $v \in V$
singular (i.e.\ $\fn^+v = 0$ and $t \cdot v = \lambda(t)v$ for all $t \in \at$
and some $\lambda \in \at^*$) such that $V = \U\g \cdot v$.
\end{definition}

Note that in fact $\U\fn^- \cdot v = V$.

\begin{proof}
The PBW theorem (the easy direction) implies that if $x_1, \dots, x_N$ is a
basis of $\g$ then the $x_1^{a_1}\cdots x_N^{a_N}$ span $\U\g$. Taking a basis
for $\g$ ordered ``$\fn^-$, then $\at$, then $\fn^+$'', we see
$\U\g = \U\fn^- \cdot \U\at \cdot \U\fn^+$, and
$\U\fn^+ \cdot v = \U\at \cdot v = \C v$ (since $\fn^+ \cdot v = 0$).
\end{proof}

\begin{remark}
If $V$ is irreducible and finite-dimensional, it is a highest weight module
(as it has a singular vector $v$, and $\U\g \cdot v$ is a nonzero submodule).
\end{remark}

\begin{proposition}
Let $V$ be a highest weight module for $\g$ (not necessarily
finite-dimensional), and let $v_\Lambda$ be a highest weight vector with
highest weight $\Lambda \in \at^*$.
\begin{enumerate}
  \item $\at$ acts diagonalisably on $V$, so
    $V = \bigoplus_{\lambda \in D(\Lambda)} V_\lambda$, where
    \[
      D(\Lambda) = \Bigl\{\Lambda - \sum k_i\alpha_i \Bigm| k_i \in \Z_{\geq 0}\Bigr\}
                 = \setc{\mu \in \at^*}{\mu \leq \Lambda}
    \]
    (the ``descent'' of $\Lambda$).
  \item $V_\Lambda = \C v_\Lambda$, and all other weight spaces are
    finite-dimensional.
  \item $V$ is irreducible if and only if all singular vectors are in
    $V_\Lambda$.
  \item $\Omega$ acts on $V$ as $|\Lambda + \rho|^2 - |\rho|^2$.
  \item If $v_\lambda \in V_\lambda$ is a singular vector, then
    $|\lambda + \rho|^2 = |\Lambda + \rho|^2$.
  \item If $\Lambda \in \R R$, then there exist only finitely many $\lambda$
    such that $V_\lambda$ contains a singular vector.
  \item $V$ contains a unique maximal proper submodule $I$. $I$ is graded by
    $\at$, $I = \bigoplus (I \cap V_\lambda)$, and $I$ is the sum of all
    proper submodules of $V$.
\end{enumerate}
\end{proposition}

\begin{remark}
The assumption in (6) is in fact unnecessary, but there isn't quite a
one-line proof in that case. Note that if $V$ is finite-dimensional, then we
in fact know that all weights lie in $\Q R$, but we are not assuming finite
dimensionality here.
\end{remark}

\begin{proof}
As $V = \U\fn^- \cdot v_\Lambda$, the expressions
$e_{-\beta_1}e_{-\beta_2}\cdots e_{-\beta_r}v_\Lambda$ span $V$ (where
$\beta_i \in R^+$ and $e_{-\beta_i} \in \g_{-\beta_i}$). The weight of such an
expression is $\Lambda - \beta_1 - \beta_2 - \dots - \beta_r$.

\begin{exercise}
Prove it!
\end{exercise}

This implies (1) immediately, and (2) upon noticing that there are finitely
many ways of writing $\lambda = \Lambda - \beta_1 - \dots - \beta_r$ for
$\beta_i \in R^+$.

(3): If $v_\lambda \in V_\lambda$ is a singular vector with
$\lambda \neq \Lambda$, then
$N = \U\g \cdot v_\lambda = \U\fn^- \cdot v_\lambda$ is a submodule of $V$ with
weights in $D(\lambda)$. Since $\lambda < \Lambda$, we have
$\Lambda \notin D(\lambda)$, whence $N$ is a proper submodule. Thus, if there
is a singular vector outside $V_\Lambda$, then $V$ is not irreducible.

Conversely, if $N \subseteq V$ is a proper submodule, then $\at N \subseteq N$,
so $N$ is graded by $\at$, and its weights lie in $D(\Lambda)$. Let
$\lambda = \Lambda - \sum k_i\alpha_i$ be a weight of $N$ with
$\alpha_i \in \Pi$ and $\sum k_i$ minimal. Since $N$ is a proper submodule,
$\lambda \neq \Lambda$, so $\sum k_i > 0$, and in that case any nonzero vector
in $N_\lambda$ is singular (since $\lambda + \alpha_i$ is not a weight of $N$
for any simple $\alpha_i$).

(7): Any proper submodule of $V$ is $\at$-graded and does not contain
$v_\Lambda$. The sum of all proper submodules still does not contain
$v_\Lambda$, so it is the maximal proper submodule.

(4): For the singular vector $v_\Lambda$ we have
$\Omega v_\Lambda = (|\Lambda + \rho|^2 - |\rho|^2)v_\Lambda$. Since
$V = \U\g \cdot v_\Lambda$ and $\Omega$ commutes with $\U\g$, we see that
$\Omega$ acts by the same constant on all of $V$.

(5) is now immediate (apply the same observation to any other singular
vector).

(6): If $V_\lambda$ contains a singular vector, then
$|\lambda + \rho|^2 = |\Lambda + \rho|^2$. This means that $\lambda$ lies in
the intersection of a sphere in $\R R$ (compact) and $D(\Lambda)$ (discrete),
so there are only finitely many such $\lambda$.
\end{proof}

\begin{definition}
Let $\Lambda \in \at^*$. A \emph{Verma module} $M(\Lambda)$ with highest
weight $\Lambda$ and highest weight vector $v_\Lambda$ is a universal module
with highest weight $\Lambda$. That is, given any other highest weight module
$W$ with highest weight $\Lambda$ and highest weight vector $w$, there exists a
unique map $M(\Lambda) \to W$ sending $v_\Lambda \mapsto w$.
\end{definition}

\begin{remark}
Verma modules were also discovered by Harish-Chandra.
\end{remark}

\begin{proposition}
Let $\Lambda \in \at^*$. Then
\begin{enumerate}
  \item there exists a unique Verma module $M(\Lambda)$;
  \item there exists a unique irreducible highest-weight module of weight
    $\Lambda$, called $L(\Lambda)$.
\end{enumerate}
\end{proposition}

\begin{proof}
Uniqueness of $M(\Lambda)$ is clear from universality. For existence, take
\[
  M(\Lambda) = \U\g \otimes_{\U\fb} \C_\Lambda,
\]
where $\fb = \fn^+ + \at$ and $\C_\Lambda$ is the $\fb$-module on which
$\fn^+ \cdot v = 0$ and $t \cdot v = \Lambda(t)v$.

\begin{remark}
You should think of this as the induced representation.
\end{remark}

More concretely,
\[
  M(\Lambda) = \U\g/\bigl(\text{left ideal generated by } u - \Lambda(u) \text{ for all } u \in \U\fb\bigr)
\]
(where $\Lambda$ is extended to an algebra homomorphism $\U\fb \to \C$ that
kills $\fn^+$). That is, $M(\Lambda)$ is the module generated by $\g$ acting
on $\mathbf{1}$ with relations
$\{\fn^+\mathbf{1} = 0,\ t\mathbf{1} = \Lambda(t)\mathbf{1} \text{ for all } t \in \at\} = J_\Lambda$,
and only the relations these imply.

If $W$ is any other highest weight module of highest weight $\Lambda$, then
$W = \U\g/J$ for some $J \supseteq J_\Lambda$, i.e.\
$M(\Lambda) \twoheadrightarrow W$.

In particular, an irreducible highest-weight module must be of the form
$M(\Lambda)/I(\Lambda)$, where $I(\Lambda)$ is a maximal proper submodule of
$M(\Lambda)$. Since we showed that there exists a unique maximal proper
submodule, $L(\Lambda)$ is unique.
\end{proof}

\begin{proposition}
Let $R^+ = \{\beta_1, \dots, \beta_r\}$. Then the vectors
$e_{-\beta_1}^{k_1}\cdots e_{-\beta_r}^{k_r}v_\Lambda$, $k_i \in \Z_{\geq 0}$,
form a basis for $M(\Lambda)$.
\end{proposition}

\begin{proof}
Immediate from the difficult direction of the PBW theorem.
\end{proof}

\begin{corollary}
Any irreducible finite-dimensional $\g$-module is of the form $L(\Lambda)$ for
some $\Lambda \in P^+$.
\end{corollary}

\begin{proof}
We know that the highest weight must lie in $P^+$ by $\sl_2$ theory.
\end{proof}

\begin{example}
$\g = \sl_2$. The Verma module $M(\Lambda)$ looks like this:
\[
\begin{tikzpicture}[baseline=-3pt,x=1.5cm]
  \foreach \i/\lab in {0/{v_\Lambda},1/{Fv_\Lambda},2/{F^2v_\Lambda}} {
    \node[dot] (n\i) at (\i,0) {};
    \node[below=2pt] at (n\i) {$\lab$};
  }
  \node (n3) at (3,0) {$\cdots$};
  \draw (n0)--(n1)--(n2)--(n3);
\end{tikzpicture}
\]
i.e.\ an infinite string.
\end{example}

\begin{exercise}\leavevmode
\begin{enumerate}
  \item Show that $M(\lambda) = L(\lambda)$ unless $\lambda \in \Z_{\geq 0}$.
  \item If $\lambda \in \Z_{\geq 0}$, show that $M(\lambda)$ contains a unique
    nonzero proper submodule, spanned by $F^{\lambda+1}v_\lambda$,
    $F^{\lambda+2}v_\lambda, \dots$ (which is itself a Verma module). Recover
    the finite-dimensional irreducible representations.
\end{enumerate}
\end{exercise}

\begin{figure}[!htbp]
\centering
\begin{tikzpicture}[scale=0.9,every node/.append style={font=\small}]
  \coordinate (L) at (0,0);
  \coordinate (A) at (-1.4,0);
  \coordinate (B) at (-60:1.4);
  \fill[gray!12] (L) -- ++(-4,0) -- ++(-60:3.2) -- ($(L)+(-60:3.2)$) -- cycle;
  \fill[gray!30,opacity=0.8] (A) -- ++(-2.6,0) -- ++(-60:3.2) -- ($(A)+(-60:3.2)$) -- cycle;
  \fill[gray!30,opacity=0.8] (B) -- ++(-4,0) -- ++(-60:1.8) -- ($(B)+(-60:1.8)$) -- cycle;
  \draw (L) -- ++(-4,0); \draw (L) -- ++(-60:3.2);
  \draw (A) -- ++(-60:3.2); \draw (B) -- ++(-4,0);
  \foreach \p in {L,A,B} \node[dot] at (\p) {};
  \node[above] at (L) {$\lambda$};
  \node[above] at (A) {$F_2^{k_2}\lambda$};
  \node[right=2pt] at (B) {$F_1^{k_1}\lambda$};
\end{tikzpicture}
\caption{Some of the singular vectors in a Verma module for $\sl_3$; the label
$F_i^{k_i}\lambda$ marks the weight of $F_i^{k_i}v_\lambda$. Note that the
quotient isn't $W$-invariant, so there must be others.}
\label{fig:verma-sl3}
\end{figure}

\begin{exercise}\leavevmode
\begin{enumerate}
  \item If $\g$ is an arbitrary simple Lie algebra and
    $\Lambda(H_i) \in \Z_{\geq 0}$, show that
    $F_i^{\Lambda(H_i)+1} \cdot v_\Lambda$ is a singular vector of
    $M(\Lambda)$. (There will be others! See figure~\ref{fig:verma-sl3}.)
  \item Compute $\ch M(\Lambda)$ for all Verma modules.
\end{enumerate}
\end{exercise}

\lecture{19} We will now try to show that $L(\Lambda)$ is finite-dimensional
if $\Lambda \in P^+$.

\begin{proposition}
If $\Lambda \in P^+$, $L(\Lambda)$ is integrable, i.e.\ the $E_i$ and $F_i$
act locally nilpotently.
\end{proposition}

\begin{proof}
If $V$ is any highest-weight module, $E_i$ acts locally nilpotently (as
$E_iV_\lambda \subseteq V_{\lambda+\alpha_i}$ and the weights of $V$ lie in
$D(\Lambda)$). We must show that $F_i$ acts locally nilpotently.

By the exercise, $F_i^{\Lambda(H_i)+1} \cdot v_\Lambda$ is a singular vector
(or zero). Since $L(\Lambda)$ is irreducible, it has no singular vectors
outside $\C v_\Lambda$, so $F_i^{\Lambda(H_i)+1}v_\Lambda = 0$.

\begin{exercise}\leavevmode
\begin{enumerate}
  \item $a^kb = \sum_{i=0}^k \binom{k}{i}\bigl((\ad a)^ib\bigr)a^{k-i}$.
  \item Using the above and the Serre relations $(\ad e_\alpha)^4e_\beta = 0$,
    show $F_i^Ne_{-\beta_1}\cdots e_{-\beta_r}v_\Lambda = 0$ for $N \gg 0$ by
    induction on $r$. \qedhere
\end{enumerate}
\end{exercise}
\end{proof}

\begin{remark}
In a general Kac--Moody algebra, the Serre relations are no longer
$(\ad e_\alpha)^4e_\beta = 0$, but the exponents are still bounded above, by
$1 - a_{ij}$ for the relevant $i$, $j$ in the generalised Cartan matrix.
\end{remark}

\begin{corollary}
$\dim L(\Lambda)_\mu = \dim L(\Lambda)_{w\mu}$ for all $w \in W$.
\end{corollary}

\begin{proof}
When we were proving Proposition~8.2, we only needed local nilpotence of $E_i$
and $F_i$ to conclude invariance under $w = s_{\alpha_i} = s_i$. Since $W$ is
generated by the $s_i$, we conclude invariance under all of $W$.
\end{proof}

\begin{theorem}[Cartan's theorem]
If $\g$ is finite-dimensional and $\Lambda \in P^+$, then $L(\Lambda)$ is
finite-dimensional.
\end{theorem}

\begin{remark}
This is clear pictorially: $L(\Lambda)$ is $W$-invariant and sits in the cone
$D(\Lambda)$. However, it's faster to work through it formally.
\end{remark}

\begin{proof}
First, we show that if $\alpha \in R^+$ then $e_{-\alpha}$ acts locally
nilpotently on $L(\Lambda)$, i.e.\ all (not just the simple) local
$\sl_2$'s act integrably. In fact, $e_{-\alpha}^nv_\Lambda = 0$ for
$n = (\Lambda, \alpha\cv) + 1$, since otherwise we have
$L(\Lambda)_{\Lambda - n\alpha} \neq 0$ and hence $s_\alpha(\Lambda - n\alpha)$
is a weight of $L(\Lambda)$. However,
\[
  s_\alpha(\Lambda - n\alpha) = s_\alpha(\Lambda) + n\alpha
    = \Lambda - (\Lambda, \alpha\cv)\alpha + \bigl((\Lambda, \alpha\cv) + 1\bigr)\alpha
    = \Lambda + \alpha > \Lambda,
\]
and so it cannot be a weight in $L(\Lambda)$, as $\Lambda$ is the highest
weight. Applying Exercise~87, we see that $e_{-\alpha}$ acts locally
nilpotently on all of $L(\Lambda)$.

Therefore,
\[
  L(\Lambda) = \U(\fn^-)v_\Lambda
    = \Span\setc{e_{-\beta_1}^{k_1}\cdots e_{-\beta_r}^{k_r}v_\Lambda}{k_i \geq 0},
    \qquad R^+ = \{\beta_1, \dots, \beta_r\},
\]
and only finitely many of these spanning vectors are nonzero, so $L(\Lambda)$
is finite-dimensional.
\end{proof}

\begin{lemma}
$s_i(R^+ \setminus \{\alpha_i\}) = R^+ \setminus \{\alpha_i\}$.
\end{lemma}

\begin{proof}
Let $\alpha \in R^+ \setminus \{\alpha_i\}$,
$\alpha = \sum_{j \neq i} k_j\alpha_j + k_i\alpha_i$ with all $k_j \geq 0$ and
some $k_j \neq 0$ for $j \neq i$ (since the only multiples of $\alpha_i$ in $R$
are $\pm\alpha_i$). Note that $s_i(\alpha)$ has the same coefficients on the
$\alpha_j$ for $j \neq i$. Since $R = R^+ \sqcup -R^+$, we see
$s_i(\alpha) \notin -R^+ \implies s_i(\alpha) \in R^+$; and
$s_i(\alpha) \neq \alpha_i$ since $s_i(\alpha_i) = -\alpha_i \neq \alpha$.
\end{proof}

\begin{lemma}
Recall $\rho = \frac12 \sum_{\alpha \in R^+} \alpha$. Then $\rho(H_i) = 1$ for
all $i$, i.e.\ $\rho = \omega_1 + \dots + \omega_l$.
\end{lemma}

\begin{proof}
Observe
\[
  s_i\rho = s_i\Bigl(\frac12\alpha_i + \frac12\sum_{\substack{\alpha \neq \alpha_i\\ \alpha \in R^+}} \alpha\Bigr)
    = -\frac12\alpha_i + \frac12\sum_{\substack{\alpha \neq \alpha_i\\ \alpha \in R^+}} \alpha
    = \rho - \alpha_i.
\]
On the other hand, $s_i\rho = \rho - (\rho, \alpha_i\cv)\alpha_i$, so
$(\rho, \alpha_i\cv) = 1$ for all $i$.
\end{proof}

\begin{lemma}[Key lemma]
Let $\Lambda \in P^+$, $\mu + \rho \in P^+$, $\mu \leq \Lambda$. If
$|\Lambda + \rho| = |\mu + \rho|$ then $\Lambda = \mu$.
\end{lemma}

\begin{remark}
This, again, is geometrically obvious from the fact that both $D(\Lambda)$ and
a sphere are convex.
\end{remark}

\begin{proof}
\[
  0 = (\Lambda + \rho, \Lambda + \rho) - (\mu + \rho, \mu + \rho)
    = \bigl(\Lambda - \mu, (\Lambda + \rho) + (\mu + \rho)\bigr).
\]
Write $\Lambda - \mu = \sum k_i\alpha_i$ with $k_i \geq 0$. Since
$\Lambda + (\mu + \rho) \in P^+$ and $(\rho, \alpha_i) > 0$ for all $i$, each
$\bigl(\alpha_i, (\Lambda + \rho) + (\mu + \rho)\bigr)$ is positive, so we
must have $k_i = 0$ for all $i$.
\end{proof}

\begin{theorem}[Weyl's complete reducibility theorem]
Let $\g$ be a semisimple Lie algebra over a field $k$ with $\charac k = 0$ and
$k = \overline{k}$. Then every finite-dimensional $\g$-module $V$ is a direct
sum of irreducibles.
\end{theorem}

\begin{proof}
Recall that $\at$ acts diagonalisably on $V$ (since $V$ is completely
reducible as an $(\sl_2)_\alpha$-module for each $\alpha$),
$V = \bigoplus V_\lambda$.

Consider $V^{\fn^+} = \setc{x \in V}{\fn^+x = 0}$. By Engel's theorem,
$V^{\fn^+} \neq 0$. Since $[\at, \fn^+] \subseteq \fn^+$, $\at$ acts on
$V^{\fn^+}$, so $V^{\fn^+} = \bigoplus_{\mu \in P} V_\mu^{\fn^+}$.

\begin{claim}
For every nonzero $v_\mu \in V_\mu^{\fn^+}$, the module $L = \U\g \cdot v_\mu$
is irreducible.
\end{claim}

\begin{proof}
$L$ is a highest weight module with highest weight $\mu$, so we must only show
that it has no other singular vectors. If $\lambda$ is the weight of a singular
vector in $L$, then $\lambda \leq \mu$, but $|\lambda + \rho| = |\mu + \rho|$
by considering the action of the Casimir. Since $V$, and therefore $L$, is
finite-dimensional, we must have $\lambda, \mu \in P^+$, so in particular
$\lambda + \rho \in P^+$. By the key lemma, $\lambda = \mu$.
\end{proof}

Consequently, $V' = \U\g \cdot V^{\fn^+}$ is completely reducible: if
$\{v_1, \dots, v_r\}$ is a weight basis for $V^{\fn^+}$ with weights
$\lambda_1, \dots, \lambda_r$, then $V' = L(\lambda_1) + \dots + L(\lambda_r)$
is a sum of irreducible submodules, and hence a direct sum of some of them.
It remains to show that $N = V/V' = 0$.

Suppose not. $N$ is finite-dimensional, so $N^{\fn^+} \neq 0$. Let
$v_\lambda \in N_\lambda^{\fn^+}$ be a singular vector, so $\lambda \in P^+$.
Lift $v_\lambda$ to $\overline{v_\lambda} \in V_\lambda$. Since
$v_\lambda \neq 0$ in $N$, $\overline{v_\lambda} \notin V^{\fn^+}$, so we must
have $E_i\overline{v_\lambda} \neq 0 \in V_{\lambda+\alpha_i}$ for some $i$.

Note that $E_iv_\lambda = 0$, so $E_i\overline{v_\lambda} \in V'$.
Consequently, the Casimir $\Omega$ acts on $E_i\overline{v_\lambda}$ by
$|\overline{\lambda} + \rho|^2 - |\rho|^2$ for some
$\overline{\lambda} \in \{\lambda_1, \dots, \lambda_r\}$ (looking at a
summand $L(\overline\lambda)$ of $V'$ in which $E_i\overline{v_\lambda}$ has a
nonzero component). On the other hand, since $v_\lambda$ was singular in $N$,
$\Omega$ acts on $v_\lambda$ by $|\lambda + \rho|^2 - |\rho|^2$. Since $\Omega$
commutes with $E_i$, we find $|\overline{\lambda} + \rho| = |\lambda + \rho|$.
Moreover, $\lambda + \alpha_i$ is a weight in $L(\overline{\lambda})$, so
$\lambda + \alpha_i = \overline{\lambda} - \sum k_j\alpha_j$, i.e.\
$\lambda < \overline{\lambda}$. This, however, contradicts the key lemma.
\end{proof}

\begin{remark}
The usual proof (in particular, Weyl's own proof) of this theorem involves
lifting the integrable action to the group, and using the maximal compact
subgroup (for which we have complete reducibility because it has positive,
finite measure). Somehow, the key lemma is the equivalent of a positive-definite
metric.
\end{remark}

\begin{lemma}
Let $\Lambda \in \at^*$ and $M(\Lambda)$ the Verma module. Then
\[
  \ch M(\Lambda) = \frac{e^\Lambda}{\prod_{\alpha \in R^+} (1 - e^{-\alpha})}.
\]
\end{lemma}

\lecture{20}

\begin{proof}
Let $R^+ = \{\beta_1, \dots, \beta_r\}$. By PBW, the basis for $M(\Lambda)$ is
$\setc{e_{-\beta_1}^{k_1}\cdots e_{-\beta_r}^{k_r}v_\Lambda}{k_i \in \Z_{\geq 0}}$,
and the weight of such an element is $\Lambda - \sum k_i\beta_i$. Therefore,
$\dim M(\Lambda)_{\Lambda - \beta}$ is the number of ways of writing $\beta$
as $\sum k_i\beta_i$ over $\beta_i \in R^+$. This is precisely the coefficient
of $e^{-\beta}$ in $\prod_{\alpha \in R^+} (1 - e^{-\alpha})^{-1}$.
\end{proof}

Write $\Delta = \prod_{\alpha \in R^+} (1 - e^{-\alpha})$. We have shown that
$\ch M(\Lambda) = e^\Lambda/\Delta$.

\begin{lemma}
For all $w \in W$, $w(e^\rho\Delta) = \det w \cdot e^\rho\Delta$. Here,
$\det\colon W \to \{\pm 1\}$ gives the determinant of $w$ acting on $\at^*$.
\end{lemma}

\begin{proof}
Since $W$ is generated by simple reflections, it suffices to check that
$s_i(e^\rho\Delta) = -e^\rho\Delta$. Indeed,
\[
  s_i(e^\rho\Delta)
  = s_i\Bigl(e^\rho(1 - e^{-\alpha_i})\prod_{\substack{\alpha \neq \alpha_i\\ \alpha \in R^+}} (1 - e^{-\alpha})\Bigr)
  = e^{\rho - \alpha_i}(1 - e^{\alpha_i})\prod_{\substack{\alpha \neq \alpha_i\\ \alpha \in R^+}} (1 - e^{-\alpha})
  = -e^\rho\Delta. \qedhere
\]
\end{proof}

\begin{lemma}\leavevmode
\begin{enumerate}
  \item For any highest weight module $V(\Lambda)$ with highest weight
    $\Lambda$ there exist coefficients $a_\lambda \geq 0$, $\lambda \leq \Lambda$,
    such that
    \[
      \ch V(\Lambda) = \sum_{\substack{\lambda \leq \Lambda\\ |\lambda + \rho| = |\Lambda + \rho|}} a_\lambda \ch L(\lambda),
    \]
    with $a_\Lambda = 1$.
  \item There exist coefficients $b_\lambda \in \Z$ with $b_\Lambda = 1$ such
    that
    \[
      \ch L(\Lambda) = \sum_{\substack{\lambda \leq \Lambda\\ |\lambda + \rho| = |\Lambda + \rho|}} b_\lambda \ch M(\lambda).
    \]
\end{enumerate}
\end{lemma}

We write
$B(\Lambda) = \setc{\lambda \leq \Lambda}{|\lambda + \rho| = |\Lambda + \rho|}$.
Recall that $B(\Lambda)$ is a finite set.

\begin{remark}
We have shown that $B(\Lambda)$ is finite when $\Lambda \in \R R$, but it is
true in general.
\end{remark}

\begin{proof}
First, note that (1) implies (2). Indeed, totally order $B(\Lambda)$ so that
$\lambda_i \leq \lambda_j \implies i \leq j$. Applying (1) to the Verma
modules gives a system of equations relating the $\ch M(\lambda)$ and the
$\ch L(\lambda)$ which is upper-triangular with $1$'s on the diagonal.
Inverting this system gives (2). (Note in particular that the coefficients in
(2) may be negative.)

We now prove (1). Recall that the weight spaces of a highest weight module are
finite-dimensional. We induct on $\sum_{\mu \in B(\Lambda)} \dim V(\Lambda)_\mu$.

If $V(\Lambda)$ is irreducible, we clearly have (1). Otherwise, there exists
$\mu \in B(\Lambda)$, $\mu \neq \Lambda$, with a singular vector
$v_\mu \in V(\Lambda)_\mu$. Pick such a $\mu$ with the largest height of
$\Lambda - \mu$ (recall the height of $\sum k_i\alpha_i$ is $\sum k_i$). Then
$\U\g \cdot v_\mu \subseteq V(\Lambda)$ has no singular vectors other than
those in its highest weight space, and is therefore irreducible, $L(\mu)$.
Setting $\overline{V(\Lambda)} = V(\Lambda)/L(\mu)$, we see that
$\overline{V(\Lambda)}$ is a highest-weight module with a smaller value of
$\sum_{\mu \in B(\Lambda)} \dim \overline{V(\Lambda)}_\mu$, and
$\ch V(\Lambda) = \ch \overline{V(\Lambda)} + \ch L(\mu)$.
\end{proof}

We will now compute $\ch L(\Lambda)$ for $\Lambda \in P^+$. We know
\[
  \ch L(\Lambda) = \sum_{\lambda \in B(\Lambda)} b_\lambda \ch M(\lambda)
    = \sum_{\lambda \in B(\Lambda)} b_\lambda e^\lambda/\Delta.
\]
Now, $w(\ch L(\Lambda)) = \ch L(\Lambda)$ for all $w \in W$, and
$w(\Delta e^\rho) = \det w \cdot \Delta e^\rho$. Therefore,
\[
  e^\rho\Delta \ch L(\Lambda) = \sum_{\lambda \in B(\Lambda)} b_\lambda e^{\lambda+\rho}
\]
is $W$-antiinvariant. We can therefore rewrite
\[
  \sum_{\lambda \in B(\Lambda)} b_\lambda e^{\lambda+\rho}
  = \sum_{\text{orbits of } W \text{ on } B(\Lambda) + \rho} b_\lambda
    \sum_{w \in W} \det w \cdot e^{w(\lambda+\rho)}.
\]
Now, if $\lambda \in \R R$ (which is true, since $\Lambda \in P^+$), then
$W(\lambda + \rho)$ intersects
$\setc{x \in \R R}{(x, \alpha_i\cv) \geq 0 \ \forall i}$ in exactly one point
(this set is a fundamental domain for the $W$-action, and we won't have
boundary problems because $(\rho, \alpha_i\cv) = 1 > 0$). Therefore, we can
take a dominant weight as the orbit representative. Note, however, that there
is only one $\lambda \in B(\Lambda)$ with $\lambda + \rho$ dominant, namely
$\Lambda$! (This was the Key Lemma~10.10.) Therefore, we derive

\begin{theorem}[Weyl character formula]
\[
  \boxed{\ch L(\Lambda)
    = \frac{\sum_{w \in W} \det w \cdot e^{w(\Lambda+\rho)}}{e^\rho\prod_{\alpha \in R^+} (1 - e^{-\alpha})}
    = \sum_{w \in W} \det w \cdot \ch M\bigl(w(\Lambda + \rho) - \rho\bigr).}
\]
\end{theorem}

\begin{example}
Let $\g = \sl_2$, and write $z = e^{\alpha/2}$. Then $\C[P] = \C[z, z^{-1}]$
and $e^\rho = z$. We have
\[
  \ch L(m\alpha/2) = (z^{m+1} - z^{-(m+1)})/(z - z^{-1})
\]
as before.
\end{example}

\begin{remark}
The second expression for the character of $L(\Lambda)$ looks like an Euler
characteristic. In fact, it is: there is a complex of Verma modules whose
cohomology is the irreducible module. This is known as the
``Bernstein--Gelfand--Gelfand'' resolution.

The first expression comes from something like a cohomology of coherent line
bundles, and is variously known as Euler--Grothendieck, Atiyah--Bott, and
Atiyah--Singer.
\end{remark}

\begin{corollary}
Since $L(0) = \C$, we must have $\ch L(0) = 1$, so
\[
  e^\rho\prod_{\alpha \in R^+} (1 - e^{-\alpha}) = \sum_{w \in W} \det w \cdot e^{w\rho}.
\]
This is known as the \emph{Weyl denominator identity}.
\end{corollary}

\begin{exercise}
Let $\g = \sl_n$. Show that the Weyl denominator identity is equivalent to the
Vandermonde determinant
\[
  \det \begin{pmatrix}
    1 & 1 & \cdots & 1\\
    z_1 & z_2 & \cdots & z_n\\
    \vdots & \vdots & \ddots & \vdots\\
    z_1^{n-1} & z_2^{n-1} & \cdots & z_n^{n-1}
  \end{pmatrix} = \prod_{i<j} (z_j - z_i),
\]
where we write $z_i = e^{e_i}$.
\end{exercise}

\begin{corollary}[Weyl dimension formula]
\[
  \dim L(\Lambda) = \prod_{\alpha \in R^+} \frac{(\alpha, \Lambda + \rho)}{(\alpha, \rho)}.
\]
\end{corollary}

\begin{example}
$\g = \sl_3$ of type $A_2$, $R^+ = \{\alpha, \beta, \alpha + \beta\}$ with
$\rho = \alpha + \beta = \omega_1 + \omega_2$. Let
$\Lambda = m_1\omega_1 + m_2\omega_2$; then
\[
  \begin{array}{c|c|c|c}
    & \alpha & \beta & \alpha + \beta\\ \hline
    (\,\cdot\,, \Lambda + \rho) & m_1 + 1 & m_2 + 1 & m_1 + m_2 + 2\\
    (\,\cdot\,, \rho) & 1 & 1 & 2
  \end{array}
\]
Therefore, the dimension of $L(\Lambda)$ is
$\frac12(m_1 + 1)(m_2 + 1)(m_1 + m_2 + 2)$.
\end{example}

``Just to check that you too certainly know the dimensions of the irreducible
representations, let me state the obvious exercise.''

\begin{exercise}
Compute the dimensions of all the finite-dimensional irreducible
representations of $B_2$ and $G_2$.
\end{exercise}

\lecture{21}

\begin{remark}
Let $w \in W$ be written as $w = s_{i_1}s_{i_2}\cdots s_{i_r}$ where the
$s_{i_k}$ are simple reflections. Then $\det w = (-1)^r$. The minimal $r$ such
that $w$ can be written in this form is called the \emph{length} of $w$,
denoted $l(w)$. The \emph{Monoid Lemma} (also known as Matsumoto's theorem)
asserts that you can get from one minimal-length expression for $w$ to
another by repeatedly applying the braid relations.
\end{remark}

\begin{exercise}
$l(w) = \#\{-R^+ \cap w^{-1}R^+\} = l(w^{-1})$.
\end{exercise}

We now go back to our proof of the Weyl dimension formula.

\begin{proof}
We know $\ch L(\Lambda) = \sum \dim L(\Lambda)_\lambda e^\lambda \in \C[P]$. We
would like to set $e^\lambda \mapsto 1$, but that leads to a $0/0$ expression
in the Weyl character formula, which is hard to evaluate. Therefore, we will
start by defining the homomorphism $F_\mu\colon \C[P] \to \C(q)$,
$e^\lambda \mapsto q^{-(\lambda, \mu)}$. We would like to evaluate
$F_0(\ch L(\Lambda))$, since $F_0(e^\mu) = 1$ for all $\mu$.

We apply $F_\mu$ to the Weyl denominator identity:
\[
  q^{-(\rho, \mu)}\prod_{\alpha \in R^+} (1 - q^{(\alpha, \mu)})
  = \sum_{w \in W} \det w \, q^{-(w\rho, \mu)}
  = \sum_{w \in W} \det w \, q^{-(\rho, w^{-1}\mu)},
\]
since $\det w = \det w^{-1}$ and $(x, wy) = (w^{-1}x, y)$ (i.e.\ the Weyl group
is a subgroup of the orthogonal group of the inner product).

We now apply $F_\mu$ to the Weyl character formula:
\[
  F_\mu(\ch L(\Lambda))
  = \frac{\sum_{w \in W} \det w \, q^{-(w(\Lambda+\rho), \mu)}}{q^{-(\rho, \mu)}\prod_{\alpha \in R^+} (1 - q^{(\alpha, \mu)})},
\]
provided $(\alpha, \mu) \neq 0$ for all $\alpha \in R^+$.

Take $\mu = \rho$ (recall that $(\rho, \alpha_i) > 0$ for all simple roots
$\alpha_i$, so $(\rho, \alpha) > 0$ for all $\alpha \in R^+$). Then
\[
  F_\rho(\ch L(\Lambda)) = \sum \dim L(\Lambda)_\lambda q^{-(\lambda, \rho)}
  = \frac{q^{-(\rho, \Lambda+\rho)}\prod_{\alpha \in R^+} (1 - q^{(\alpha, \Lambda+\rho)})}
         {q^{-(\rho, \rho)}\prod_{\alpha \in R^+} (1 - q^{(\alpha, \rho)})},
\]
where we used our expression for the Weyl denominator identity, with
$\mu = \Lambda + \rho$, to rewrite the numerator.

Setting $q = 1$ and applying L'H\^opital's rule,
\[
  \dim L(\Lambda) = \prod_{\alpha \in R^+} \frac{(\Lambda + \rho, \alpha)}{(\rho, \alpha)}. \qedhere
\]
\end{proof}

\begin{remark}
``You are now in a position to answer any question you choose to answer, if
necessary by brute force.'' For example, let us decompose
$L(\lambda) \otimes L(\mu) = \sum m_{\lambda\mu}^\nu L(\nu)$; we want to
compute the \emph{Littlewood--Richardson coefficients} $m_{\lambda\mu}^\nu$
(recall that we had the Clebsch--Gordan rule for them in $\sl_2$).

Define $\overline{\phantom{x}}\colon \Z[P] \to \Z[P]$, $e^\lambda \mapsto e^{-\lambda}$,
and $\CT\colon \Z[P] \to \Z$ with $e^\lambda \mapsto 0$ unless $\lambda = 0$,
and $e^0 \mapsto 1$. Define $(\,,\,)\colon \Z[P] \times \Z[P] \to \Z$ by
$(f, g) = \frac{1}{|W|}\CT(f\overline{g}\Delta\overline{\Delta})$, where
$\Delta = \prod_{\alpha \in R^+} (1 - e^{-\alpha})$.
\end{remark}

\begin{claim}
Let $\chi_\lambda = \ch L(\lambda)$, $\chi_\mu = \ch L(\mu)$; then
$(\chi_\lambda, \chi_\mu) = \delta_{\lambda\mu}$. (If our Lie algebra has type
$A$, then the $\chi_\lambda$ are the Schur functions, and this is their
orthogonality.)
\end{claim}

Given this claim, $m_{\lambda\mu}^\nu = (\chi_\lambda\chi_\mu, \chi_\nu)$ is
algorithmically computable. It was a great surprise to everyone when it was
discovered that there is a simpler and more beautiful way of doing this.

\begin{proof}[Proof of claim]
By the Weyl character formula, $\chi_\lambda\Delta = \sum_{w \in W} \det w\,
e^{w(\lambda+\rho) - \rho}$. Therefore
\[
  (\chi_\lambda, \chi_\mu) = \frac{1}{|W|}\CT\Bigl(\sum_{x, w \in W}
    e^{w(\lambda+\rho) - \rho}\,\overline{e^{x(\mu+\rho) - \rho}}\det(wx)\Bigr).
\]
However, $\CT(e^{w(\lambda+\rho) - x(\mu+\rho)}) = \delta_{wx}\delta_{\mu\lambda}$,
since for $\lambda, \mu \in P^+$ we have
$x^{-1}w(\lambda + \rho) = \mu + \rho \implies x = w,\ \lambda = \mu$.
\end{proof}
